% =============================================================================
%  vsim_cosmology_qa.tex
%  Cosmology Q&A — Open Questions and Current Understanding
%  Internal working document · v5.2.0-internal context
% =============================================================================
\documentclass[11pt,a4paper]{article}
\usepackage{vsim_common}

% Q&A specific layout
\usepackage{mdframed}

% Answer space rule
\newcommand{\answerspace}[1][4cm]{%
  \vspace{4pt}%
  \noindent\textcolor{gray!40}{\rule{\linewidth}{0.3pt}}\\[-2pt]%
  \vspace{#1}%
  \noindent\textcolor{gray!40}{\rule{\linewidth}{0.3pt}}%
  \vspace{6pt}%
}

% Partial answer (some text, then space)
\newcommand{\partialanswer}[2][3cm]{%
  \vspace{4pt}%
  \noindent\textcolor{gray!40}{\rule{\linewidth}{0.3pt}}\\[4pt]%
  {\small\color{black!70} #2}%
  \vspace{#1}%
  \noindent\textcolor{gray!40}{\rule{\linewidth}{0.3pt}}%
  \vspace{6pt}%
}

% Question counter
\newcounter{qnum}[section]
\newcommand{\Q}[1]{%
  \stepcounter{qnum}%
  \vspace{10pt}%
  \noindent{\sffamily\bfseries Q\thesection.\theqnum\quad #1}%
  \vspace{4pt}\par%
}

% Status tag for questions
\newcommand{\qstatus}[1]{\hfill\stbd}
\newcommand{\qopen}{\hfill\snew}
\newcommand{\qsettled}{\hfill\spass}
\newcommand{\qscope}{\hfill\sna}  % out of scope

% Note box (used for context/scope notes)
\newenvironment{scopenote}{%
  \begin{warnbox}\small\itshape%
}{%
  \end{warnbox}%
}

\hypersetup{pdftitle={Cosmology Q\&A},pdfsubject={VSIM Internal context document}}

% =============================================================================
\begin{document}
% =============================================================================

\thispagestyle{empty}

\begin{center}
  \eyebrow{VSIM Internal \textperiodcentered\ Working Document}\\[8pt]
  {\LARGE\sffamily\bfseries Cosmology Q\&A}\\[6pt]
  {\normalsize\sffamily
    Current Understanding \quad\textperiodcentered\quad
    Open Questions \quad\textperiodcentered\quad
    Theoretical Scope Notes}
\end{center}

\vspace{10pt}
\noindent
\begin{tabular}{@{}L{2.8cm} L{9cm}@{}}
  \code{Doc ID}    & QA-COSM-2026-01 \\
  \code{Status}    & Working draft --- nearly blank \\
  \code{Audience}  & Internal \\
  \code{Context}   & Background reference for identity/scale framework \\
  \code{Date}      & 2026-06-01 \\
\end{tabular}

\vspace{10pt}
\begin{infobox}
  \small\sffamily
  \textbf{How to use this document.}\quad
  Questions are grouped by topic. Each has a status tag:
  \spass\ = settled consensus \quad
  \stbd\ = active research \quad
  \snew\ = open / speculative \quad
  \sna\ = out of scope for this project.\\[4pt]
  Answer spaces are intentionally blank or partially filled.
  Fill them in as understanding develops.
  Scope notes flag where a question does or does not bear on the VSIM identity framework.
\end{infobox}

\noindent\rule{\linewidth}{2pt}\vspace{6pt}
\tableofcontents
\newpage

% =============================================================================
\section{Origins and Structure of the Universe}
% =============================================================================

\Q{What does the Big Bang Theory propose?}\qsettled

\partialanswer[2.5cm]{%
  The leading explanation for the origin of the universe. Proposes that approximately
  \textbf{13.8 billion years ago} the universe began as an infinitely hot and dense
  singularity and has been expanding and cooling ever since. Supported by the Hubble
  expansion, CMB radiation, and Big Bang nucleosynthesis predictions.
}

\Q{Is the universe still expanding?}\qsettled

\partialanswer[2.5cm]{%
  Yes. Observations of distant galaxies show recession velocities proportional to
  distance (Hubble's law). The rate of expansion is described by the Hubble constant
  $H_0 \approx 67$--$74\ \text{km\,s}^{-1}\text{Mpc}^{-1}$ (current measurement tension unresolved).
}

\Q{What is the cosmic microwave background radiation (CMB)?}\qsettled

\partialanswer[2.5cm]{%
  Residual thermal radiation from approximately 380,000 years after the Big Bang,
  when the universe cooled enough for protons and electrons to combine into neutral
  hydrogen (recombination epoch). The CMB fills the observable universe uniformly
  at $T \approx 2.725\ \text{K}$ and carries an imprint of density fluctuations that
  seeded large-scale structure.
}

\Q{How does the CMB relate to the identity/scale framework?}\qopen

\begin{scopenote}
  Peripheral scope. The CMB encodes information about the early universe's state
  that cannot be recovered at later scales --- a physical analogue of formation memory
  ($M_f$) and projection loss. Worth noting conceptually; not a current modelling target.
\end{scopenote}
\answerspace[3cm]

% =============================================================================
\section{Dark Energy}
% =============================================================================

\Q{What is dark energy?}\stbd

\partialanswer[3cm]{%
  A mysterious component responsible for the \textbf{accelerated expansion} of the
  universe, first confirmed by Type Ia supernova observations in 1998.
  Makes up approximately \textbf{68\%} of the universe's total energy density.
  Leading candidates: cosmological constant $\Lambda$ (vacuum energy),
  quintessence (dynamic scalar field), or modified gravity. Nature unknown.
}

\Q{Why does dark energy cause acceleration rather than deceleration?}\stbd

\answerspace[4.5cm]

\Q{Could dark energy change over time?}\snew

\answerspace[4.5cm]

\Q{What is the equation of state parameter $w$ for dark energy, and what does it constrain?}\stbd

\partialanswer[3cm]{%
  $w = P / \rho c^2$ where $P$ is pressure and $\rho$ the energy density.
  For a cosmological constant, $w = -1$. Current observations constrain
  $w \approx -1.0 \pm 0.05$. If $w < -1$ (phantom energy), the universe may end
  in a ``Big Rip.''
}

% =============================================================================
\section{Dark Matter}
% =============================================================================

\Q{What is dark matter?}\stbd

\partialanswer[3cm]{%
  An invisible, non-luminous form of matter inferred from gravitational effects:
  galaxy rotation curves, gravitational lensing, galaxy cluster dynamics, and
  large-scale structure formation. Makes up approximately \textbf{27\%} of the
  universe's mass-energy budget. Has not been directly detected. Leading candidates:
  WIMPs, axions, sterile neutrinos, primordial black holes.
}

\Q{How is dark matter distributed in galaxies?}\stbd

\answerspace[4cm]

\Q{What experiments are currently searching for dark matter directly?}\stbd

\answerspace[4cm]

\Q{Could dark matter and dark energy be related phenomena?}\snew

\answerspace[4cm]

\Q{Does dark matter interact with ordinary matter via any force other than gravity?}\stbd

\answerspace[4cm]

% =============================================================================
\section{Open Problems}
% =============================================================================

\Q{What caused the matter-antimatter asymmetry?}\snew

\partialanswer[3cm]{%
  The universe is observed to be dominated by matter. The Standard Model predicts
  nearly equal amounts of matter and antimatter should have been produced at the
  Big Bang --- yet almost all antimatter has annihilated, leaving a residual matter
  excess of roughly $10^{-9}$ per photon. The source of this CP violation beyond
  the Standard Model is unknown. Related to baryogenesis and leptogenesis models.
}

\Q{What is the fate of the universe?}\snew

\partialanswer[3.5cm]{%
  Depends on the nature of dark energy and $w$:\\[4pt]
  \begin{tabular}{@{}ll@{}}
    $w = -1$ (constant $\Lambda$) & Continued accelerated expansion; heat death \\
    $w > -1$ (quintessence)       & Expansion decelerates eventually \\
    $w < -1$ (phantom energy)     & ``Big Rip'' --- all structures torn apart \\
    $\Omega_\text{total} > 1$     & ``Big Crunch'' --- eventual re-collapse \\
  \end{tabular}
}

\Q{How did gravity shape the universe during its first fraction of a second?}\snew

\partialanswer[3cm]{%
  During the Planck epoch ($t < 10^{-43}\ \text{s}$), quantum gravitational effects
  dominated. At the GUT epoch ($t \sim 10^{-36}\ \text{s}$), inflation likely occurred ---
  exponential expansion driven by a scalar inflaton field --- smoothing the universe
  and amplifying quantum fluctuations into the density seeds observed in the CMB.
  The precise mechanism is unknown; quantum gravity is not yet unified with the
  Standard Model.
}

\Q{Did gravitational particles (gravitons) always exist?}\qscope

\begin{scopenote}
  Out of scope for this project's immediate modelling targets, but noted here for
  completeness. This is an ongoing open question in quantum gravity. The graviton
  is the hypothetical force-carrier of gravity in a quantum field theory framework;
  it has not been directly detected. Whether gravitational degrees of freedom
  existed in any form prior to the Planck epoch is model-dependent (string theory,
  loop quantum gravity, causal dynamical triangulations each give different pictures).
  Not a claim this framework needs to adjudicate.
\end{scopenote}
\answerspace[3cm]

\Q{What happened before the Big Bang?}\snew

\answerspace[5cm]

\Q{Is the observable universe all there is?}\snew

\answerspace[4cm]

% =============================================================================
\section{Framework Relevance Notes}
% =============================================================================

\Q{Which cosmological concepts map most directly onto the identity/scale framework?}\stbd

\partialanswer[4cm]{%
  Provisional mapping for internal reference only:

  \vspace{6pt}
  \begin{tabular}{@{}L{5cm} L{7.5cm}@{}}
    \textbf{Cosmological concept} & \textbf{Framework analogue} \\
    \toprule
    CMB information content       & Formation memory $M_f$ at maximum (early universe) \\
    Projection into observable    & Scale projection / coalescence residual $R_{\text{coal}}$ \\
    Matter-antimatter asymmetry   & Identity distinguishability $J_i$ at birth \\
    Dark matter (unobservable)    & Hidden activity / residual structure $R_i$ \\
    Hubble tension                & Variance in $Q_n$ across scale levels \\
    \bottomrule
  \end{tabular}
}

\Q{What additional cosmological questions should be tracked here?}\stbd

\answerspace[5cm]

% =============================================================================
\vspace{16pt}
\noindent\rule{\linewidth}{0.5pt}
{\small\sffamily\textcolor{gray}{QA-COSM-2026-01 \textperiodcentered\
  VSIM Internal \textperiodcentered\ Working draft \textperiodcentered\
  Nearly blank --- fill in as understanding develops}}
\end{document}
